\documentclass[12pt]{article}
\usepackage{amsmath, amssymb, amsthm}
\usepackage{geometry}
\usepackage{hyperref}
\usepackage{setspace}
\usepackage{enumitem}
\usepackage{physics}
\usepackage{xcolor}
\geometry{margin=0.75in}

\title{\textbf{PQBFL Protocol}}
\date{}

\begin{document}
\maketitle

\section*{Phase 1: Initial State Generation \& Registration}

\subsection*{1. Server (S): ``Bob''}

The server, "Bob," is the central aggregator for the Federated Learning task. Its first step is to generate all the necessary cryptographic keys and then publish a "project" to the blockchain, committing to its public keys in the process.

\subsubsection*{1.1. State Component 1: KEM Key Pair $(kpk_b, ksk_b)$}

\textbf{Description:}
This is the \textbf{Post-Quantum Key Encapsulation Mechanism (KEM)} key pair. It is the "quantum-resistant" component of the hybrid key exchange. The public key, $kpk_b$, will be used by participants (like Alice) to create and "encapsulate" a shared secret that only Bob can "decapsulate" using his secret key, $ksk_b$. The protocol uses \textbf{CRYSTALS-Kyber}.

\textbf{Kyber.Gen:}
The command $(kpk_b, ksk_b) \leftarrow \text{Kyber.KeyGen()}$ summarizes a complex process based on the \textbf{Module-Learning With Errors (M-LWE)} problem.

\begin{enumerate}
    \item \textbf{Seed Generation:} Bob first generates a 32-byte seed, $\rho$.
    \item \textbf{Matrix Generation ($A$):} This seed $\rho$ is fed into an eXtendable-Output Function (XOF) like \textbf{SHAKE-128} to deterministically generate the public matrix $A$. This way, $A$ doesn't need to be stored; it can be regenerated from $\rho$.
    \[
    A \leftarrow \text{SHAKE-128}(\rho)
    \]
    \item \textbf{Secret Sampling ($s, e$):} Bob securely generates two vectors of "small" polynomials, $s$ (the core secret key) and $e$ (the "error" or "noise"). These are sampled from a \textbf{Centered Binomial Distribution (CBD)}.
    \item \textbf{Public Vector Calculation ($t$):} Bob computes the public vector $t$ using the main M-LWE equation. All arithmetic is done in a polynomial ring $R_q = \mathbb{Z}_q[X] / (X^{256} + 1)$.
    \[
    t = A s + e
    \]
    \item \textbf{Final Keys:} The public and secret keys are then defined:
    \begin{itemize}
        \item \textbf{Public Key ($kpk_b$):} This is the pair $(t, \rho)$. $t$ is the public vector, and $\rho$ is the seed to recreate $A$.
        \item \textbf{Secret Key ($ksk_b$):} This is the secret vector $s$.
    \end{itemize}
\end{enumerate}

\textbf{Time Complexity:}
The most expensive operation is the matrix-vector multiplication $A s$.
\begin{itemize}
    \item This involves $k^2$ polynomial multiplications (where $k=3$ for Kyber-768).
    \item Each polynomial multiplication is of degree $n=255$.
    \item Using the \textbf{Number-Theoretic Transform (NTT)}, each multiplication is $O(n \log n)$ instead of $O(n^2)$.
    \item The total complexity is $O(k^2 \cdot n \log n)$.
    \item Since $k$ and $n$ are fixed by the Kyber security level (e.g., $k=3, n=256$), the operation has a fixed, constant runtime. Therefore, the complexity is \textbf{$O(1)$}.
\end{itemize}

\subsubsection*{1.2. State Component 2: Key Exchange Pair $(epk_b, esk_b)$}

\textbf{Description:}
This is the \textbf{classical} key exchange pair, based on \textbf{Elliptic Curve Diffie-Hellman (ECDH)}. It is used in parallel with Kyber to create the hybrid secret. This provides security in case the post-quantum algorithm (Kyber) is ever broken, as the attacker would \textit{also} need to break classical ECDH.

\textbf{ (ECDH.Gen):}
The generation follows the standard elliptic curve key pair process:
\begin{enumerate}
    \item \textbf{Select Curve:} A standard, secure curve is chosen, such as \textbf{NIST P-256} (as shown in the Annexure). This defines the field $K$, the equation, and the base point $G$.
    \item \textbf{Generate Secret ($esk_b$):} Bob generates a cryptographically secure random integer $esk_b$ (the secret key). This integer is in the range $[1, n-1]$, where $n$ is the order of the base point $G$.
    \item \textbf{Compute Public Key ($epk_b$):} Bob performs elliptic curve \textbf{scalar multiplication} (also called point multiplication) of the base point $G$ by his secret scalar $esk_b$ to get his public key $epk_b$, which is a new point on the curve.
    \[
    epk_b = esk_b \cdot G
    \]
\end{enumerate}

\textbf{Time Complexity:}
The cost is one scalar multiplication. For a fixed curve with $m$-bit keys (e.g., $m=256$), this is a fixed number of point additions and doublings. Thus, the complexity is \textbf{$O(1)$}.

\subsubsection*{1.3. State Component 3: Signature Key Pair $(pk_b, sk_b)$}

\textbf{Description:}
This is the classical \textbf{Elliptic Curve Digital Signature Algorithm (ECDSA)} key pair. This is Bob's "identity." He uses $sk_b$ to sign messages (both on-chain transactions and off-chain model data), and anyone can use $pk_b$ (which is public) to verify that the message came from him.

\textbf{ (ECDSA.Gen):}
The key generation math is \textbf{identical to ECDH key generation}.
\begin{enumerate}
    \item \textbf{Select Curve:} A standard curve is chosen (e.g., \texttt{secp256k1}, used by Bitcoin and Ethereum, or P-256).
    \item \textbf{Generate Secret ($sk_b$):} Bob generates a cryptographically secure random integer $sk_b$ (the secret signing key).
    \item \textbf{Compute Public Key ($pk_b$):} Bob computes the corresponding public key $pk_b$ via scalar multiplication.
    \[
    pk_b = sk_b \cdot G
    \]
    (Note: $sk_b$ and $esk_b$ are generated independently and are different secret numbers.)
\end{enumerate}

\textbf{Time Complexity:}
Same as ECDH key generation. It's one scalar multiplication, which is \textbf{$O(1)$} for a fixed curve.

\subsubsection*{1.4. Blockchain Registration (On-Chain)}

\textbf{Description:}
After generating all his keys, Bob "publishes" his project to the blockchain. He \textbf{does not} post his public keys directly, as this would be very expensive (high gas cost). Instead, he posts a \textbf{cryptographic commitment} (a hash) of his keys. This permanently and immutably links his blockchain address to the keys he will use.

\textbf{:}
\begin{enumerate}
    \item \textbf{Hashing ($H$):} Bob uses a standard cryptographic hash function (e.g., SHA3-256) to create his commitments.
    \[
    h_{keys} \leftarrow H(kpk_b \parallel epk_b)
    \]
    Here, $\parallel$ denotes the concatenation of the byte representations of the keys.
    \[
    h_{M0} \leftarrow H(M^0)
    \]
    This is a hash of the \textit{initial} global model (e.g., a randomly initialized neural network).
    \item \textbf{Transaction ($Tx_r$):} He builds a transaction to send to the smart contract:
    \[
    Tx_r \leftarrow \{h(kpk_b \parallel epk_b), n, h(M^0), id_p\}
    \]
    \begin{itemize}
        \item $h(...)$: The commitment to his public keys.
        \item $n$: The number of participants for the project.
        \item $h(M^0)$: The hash of the starting model.
        \item $id_p$: A unique project identifier.
    \end{itemize}
    \item \textbf{Smart Contract Call:} Bob signs and sends this transaction, calling the \texttt{RegisterProject} function on the smart contract.
    \[
    BC \leftarrow \text{RegisterProject}(Tx_r)
    \]
    The blockchain's consensus mechanism validates this transaction and, if successful, adds it to the global ledger $L$.
\end{enumerate}

\textbf{Time Complexity:}
\begin{itemize}
    \item \textbf{Local Computation:} Bob performs two hash operations. The cost of a hash function is linear in the size of its input, $O(L)$. Since the keys and initial model are of a fixed (or bounded) size, the hashing cost is \textbf{$O(1)$}.
    \item \textbf{On-Chain Interaction:} This is not a computational complexity but a \textbf{network latency}. The time is determined by the blockchain's block time (e.g., $\sim$15 seconds).
\end{itemize}

\subsection*{2. Participant (P): ``Alice''}

Alice is a participant (a local trainer) who wishes to join Bob's project. She must also generate her own keys and register on-chain.

\subsubsection*{2.1. State Component 1: Key Exchange Pair $(epk_a, esk_a)$}

\textbf{Description:}
Alice's classical \textbf{ECDH} key pair. This is \textit{her} half of the classical key exchange. She will use $esk_a$ to compute the shared secret with Bob's $epk_b$.

\textbf{ (ECDH.Gen):}
Identical process to Bob's, just with different variables:
\[
(epk_a, esk_a) \leftarrow \text{ECDH.KeyGen()}
\]
She generates a random scalar $esk_a$ and computes her public key:
\[
epk_a = esk_a \cdot G
\]

\textbf{Time Complexity:}
\textbf{$O(1)$} (one scalar multiplication on a fixed curve).

\subsubsection*{2.2. State Component 2: Signature Key Pair $(pk_a, sk_a)$}

\textbf{Description:}
Alice's \textbf{ECDSA} key pair. This is \textit{her} on-chain identity, used to sign her transactions (like \texttt{RegisterClient}) and her off-chain model updates.

\textbf{ (ECDSA.Gen):}
Identical process to Bob's:
\[
pk_a = sk_a \cdot G
\]

\textbf{Time Complexity:}
\textbf{$O(1)$} (one scalar multiplication on a fixed curve).

\subsubsection*{2.3. Blockchain Registration (On-Chain)}

\textbf{Description:}
Alice sees Bob's project $id_p$ on the blockchain. To join, she calls the \texttt{RegisterClient} function, committing to her public ECDH key.

\textbf{:}
\begin{enumerate}
    \item \textbf{Hashing ($H$):} Alice hashes her public key to create a commitment.
    \[
    h_{epk_a} \leftarrow H(epk_a)
    \]
    \item \textbf{Transaction ($Tx_r'$):} She builds her registration transaction:
    \[
    Tx_r' \leftarrow \{h(epk_a), id_p\}
    \]
    \begin{itemize}
        \item $h(epk_a)$: Her commitment to her public key.
        \item $id_p$: The ID of the project she wants to join.
    \end{itemize}
    \item \textbf{Smart Contract Call:} She signs and sends the transaction:
    \[
    BC \leftarrow \text{RegisterClient}(Tx_r')
    \]
    This adds her registration to the global ledger $L$. Now, Bob (and the public) can see that Alice has joined the project.
\end{enumerate}

\textbf{Time Complexity:}
\begin{itemize}
    \item \textbf{Local Computation:} One hash operation, \textbf{$O(1)$}.
    \item \textbf{On-Chain Interaction:} Network latency (block time).
\end{itemize}

\subsection*{3. Blockchain (BC)}

\textbf{Description:}
The Blockchain (BC) is the "public bulletin board" and trust anchor for the entire protocol. It's a \textbf{decentralized, append-only ledger} where the state $L$ (the list of all transactions) is replicated across all nodes in the network. Its purpose is to provide a single, immutable source of truth for public key commitments ($h_{keys}$, $h_{epk_a}$) and to log the \textit{hashes} of models for accountability.

\textbf{:}
The blockchain's state is its ledger $L$. The initial state is the "genesis block," or an empty set $L = \langle \rangle$. The state is updated by a \textbf{State Transition Function}, which is the \textbf{Consensus} algorithm (e.g., Proof-of-Work or Proof-of-Stake).
\[
L_{\text{new}} = \text{Consensus}(L_{\text{old}}, \text{NewBlock})
\]
This function ensures that a \texttt{NewBlock} is only appended if it contains valid transactions (e.g., all signatures are correct, no double-spending) and meets the consensus rules (e.g., a valid PoW hash).

\textbf{Time Complexity:}
\begin{itemize}
    \item \textbf{Read:} Reading data from the blockchain (e.g., Bob reading $Tx_r'$ to find Alice's key hash) is extremely fast. If indexed (like a database), it's \textbf{$O(1)$}.
    \item \textbf{Write:} Writing data (i.e., getting a transaction confirmed in a block) is not measured in $O$-notation. It is a \textbf{network latency} dictated by the consensus mechanism's \textbf{block time} (e.g., $\sim$15 seconds for Ethereum, $\sim$10 minutes for Bitcoin).
\end{itemize}

\newpage

\section*{Phase 2: Session Establishment (Off-Chain)}

After both parties have registered their identities and key commitments on the blockchain (Phase 1), they establish a secure off-chain channel. The goal of this phase is to derive a shared, hybrid (classical + post-quantum) root key \( RK_1 \). This flow is based on Figure 3 in the paper.



\subsection*{1. Server (S) Sends Keys}
The server (Bob) initiates the off-chain key exchange by sending his public keys to the participant (Alice).

\[
msg_b \leftarrow \langle kpk_b, epk_b, Tx_r, id_p \rangle \text{ }
\]
\[
\sigma_b \leftarrow \text{Sig.Sign}(sk_b, msg_b) \text{ }
\]

\begin{itemize}
    \item \textbf{\(msg_b\)}: Bob bundles his actual public keys (\(kpk_b\) for Kyber, \(epk_b\) for ECDH) into a message. He also includes his registration transaction \(Tx_r\) as proof of his on-chain identity.
    \item \textbf{\(\sigma_b\)}: He signs this message using his private ECDSA key (\(sk_b\)). This signature provides \textbf{authenticity} (proves the message is from him) and \textbf{integrity} (proves the message hasn't been altered).
    \item He sends the pair \( \langle msg_b, \sigma_b \rangle \) to Alice.
\end{itemize}
\textbf{Time Complexity:} \( O(1) \). A single ECDSA signature operation on fixed-size data.

\subsection*{2. Participant (P) Verifies \& Generates Secrets}
Alice receives the message and must rigorously verify Bob's identity before proceeding.

\[
\text{Sig.Ver}(pk_b, msg_b, \sigma_b) \stackrel{?}{=} 1 \text{ }
\]
\[
H(kpk_b \parallel epk_b) \stackrel{?}{=} h_{keys} \text{ (from on-chain } Tx_r) \text{ }
\]
\textbf{Verification:}
\begin{itemize}
    \item \textbf{Signature Verification:} Alice retrieves Bob's public signature key \(pk_b\) (from his on-chain address) and uses it to verify the signature \(\sigma_b\). This proves the message was signed by the owner of that on-chain identity.
    \item \textbf{Key Commitment Verification:} This is a critical step to prevent Man-in-the-Middle (MITM) attacks. Alice hashes the public keys \(kpk_b\) and \(epk_b\) she received in the message. She then compares this hash to the hash \(h_{keys}\) that Bob permanently committed to in his \(Tx_r\) on the blockchain. If they match, she knows these are Bob's authentic public keys .
\end{itemize}
If both checks pass, Alice computes the two shared secrets:
\[
SS_e \leftarrow \text{KE.Derive}(esk_a, epk_b) \text{ }
\]
\[
(ct, SS_k) \leftarrow \text{KEM.Encap}(kpk_b) \text{ }
\]
\textbf{Secret Generation:}
\begin{itemize}
    \item \textbf{\(SS_e\)}: The classical \textbf{ECDH} secret. Alice combines her \textit{private} key \(esk_a\) with Bob's \textit{public} key \(epk_b\) to compute the secret \(SS_e\).
    \item \textbf{\(SS_k\)}: The post-quantum \textbf{Kyber} secret. Alice uses Bob's \textit{public} KEM key \(kpk_b\) to generate a new secret \(SS_k\) and a ciphertext \(ct\) that encapsulates it. Only Bob, with his private KEM key, can "decapsulate" \(ct\) to get \(SS_k\).
\end{itemize}
\textbf{Time Complexity:} \( O(1) \). This is a sum of several constant-time operations: one signature verification, one hash, one ECDH derivation, and one KEM encapsulation.

\subsection*{3. Participant (P) Generates Root Key}
Alice now combines both secrets into the first root key.
\[
RK_1 \leftarrow KDF_A(SS_e \parallel SS_k, 0x00) \text{ }
\]
\begin{itemize}
    \item \textbf{\(KDF_A\)}: This is the \textbf{Asymmetric Key Derivation Function} (e.g., HKDF-Extract). Its job is to take the (potentially non-uniform) shared secrets and "extract" a single, cryptographically strong key.
    \item \textbf{\(SS_e \parallel SS_k\)}: This is the \textbf{hybrid} nature of the protocol. The root key is derived from both the classical and post-quantum secrets. An attacker must break \textit{both} ECDH and Kyber to discover \(RK_1\).
    \item \textbf{\(0x00\)}: A zero-byte salt, as specified in the paper's diagram.
\end{itemize}
\textbf{Time Complexity:} \( O(1) \). The KDF runs on a fixed-size input (the two concatenated secrets).

\subsection*{4. Participant (P) Sends Keys Back}
Alice must now send Bob the information he needs to compute the same secrets.
\[
msg_a \leftarrow \langle epk_a, Tx_r', id_p, ct \rangle \text{ }
\]
\[
\sigma_a \leftarrow \text{Sig.Sign}(sk_a, msg_a)
\]
\begin{itemize}
    \item \textbf{\(msg_a\)}: Alice's message contains her \textit{public} ECDH key \(epk_a\) (so Bob can compute \(SS_e\)) and the Kyber ciphertext \(ct\) (so Bob can compute \(SS_k\)).
    \item \textbf{\(\sigma_a\)}: She signs this message with her private ECDSA key \(sk_a\) to prove her identity and ensure message integrity.
    \item She sends \( \langle msg_a, \sigma_a \rangle \) back to Bob.
\end{itemize}
\textbf{Time Complexity:} \( O(1) \). A single signature operation.

\subsection*{5. Server (S) Verifies \& Generates Secrets}
Bob receives Alice's reply and performs the same verification steps.
\[
\text{Sig.Ver}(pk_a, msg_a, \sigma_a) \stackrel{?}{=} 1 \text{ }
\]
\[
H(epk_a) \stackrel?{=} h_{epk_a} \text{ (from on-chain } Tx_r') \text{ }
\]
\textbf{Verification:}
\begin{itemize}
    \item \textbf{Signature Verification:} Bob uses Alice's public key \(pk_a\) (from the blockchain) to verify her signature.
    \item \textbf{Key Commitment Verification:} Bob hashes the \(epk_a\) from the message and checks it against the hash \(h_{epk_a}\) Alice committed to in her on-chain registration \(Tx_r'\). This prevents MITM attacks.
\end{itemize}
If both checks pass, Bob derives the shared secrets:
\[
SS_e \leftarrow \text{KE.Derive}(esk_b, epk_a) \text{ }
\]
\[
SS_k \leftarrow \text{KEM.Decap}(ksk_b, ct) \text{ }
\]
\textbf{Secret Generation:}
\begin{itemize}
    \item \textbf{\(SS_e\)}: Bob computes the \textit{identical} classical secret by combining his \textit{private} key \(esk_b\) with Alice's \textit{public} key \(epk_a\).
    \item \textbf{\(SS_k\)}: Bob uses his \textit{private} KEM key \(ksk_b\) to "decapsulate" the ciphertext \(ct\), revealing the \textit{identical} post-quantum secret \(SS_k\).
\end{itemize}
\textbf{Time Complexity:} \( O(1) \). (Verification, hash, ECDH derivation, and KEM decapsulation).

\subsection*{6. Server (S) Generates Root Key}
Bob now has the same inputs as Alice and can derive the same key.
\[
RK_1 \leftarrow KDF_A(SS_e \parallel SS_k, 0x00)
\]
\begin{itemize}
    \item Bob feeds the exact same two secrets (\(SS_e\) and \(SS_k\)) into the same deterministic \(KDF_A\) function.
    \item This results in him computing a key \(RK_1\) that is \textbf{identical} to the one Alice computed in Step 3.
\end{itemize}
At this point, the session is established. Both Alice and Bob share the secret root key \(RK_1\), but this key was never transmitted over the network.

\textbf{Time Complexity:} \( O(1) \).

\newpage

\section*{Phase 3: Send and Receive Model Loop (Iterative)}

This is the main loop of the protocol, repeated for each training round \( r \). After establishing the hybrid root key \( RK_j \), this phase uses a \textbf{Symmetric Ratchet} (\( KDF_S \)) to derive a new, ephemeral key for each round. This provides \textbf{forward secrecy} (past keys can't be found from a current key) and \textbf{post-compromise security} (future keys are protected if an asymmetric ratchet is triggered) .

This flow is based on Figure 4 and Figure 5 in the paper.



\subsection*{1. Server (S) Publishes Task (On-Chain)}
Bob (the server) initiates the round by deriving this round's keys and publishing a task to the blockchain.

\[
(CK_{i,j}, K_{i,j}) \leftarrow KDF_S(CK_{i-1,j})
\]
\[
h_{Inf_b^r} \leftarrow H(Inf_b^r)
\]
\[
Tx_p \leftarrow \{r, h_{Inf_b^r}, id_p, id_t, D_t\} \text{ }
\]
\[
BC \leftarrow \text{PublishTask}(Tx_p)
\]

\begin{itemize}
    \item \textbf{Turn the Ratchet:} Bob applies the symmetric key derivation function \( KDF_S \) (e.g., HKDF-Expand) to the \textit{previous} round's chain key \( CK_{i-1,j} \). This deterministically produces two new keys:
    \begin{itemize}
        \item \( CK_{i,j} \): The new chain key, which is kept secret and used as the input for the \textit{next} round (r+1).
        \item \( K_{i,j} \): The **model key** for this round (r) only. It will be used for encryption and decryption and then discarded.
    \end{itemize}
    \item \textbf{Commitment:} He hashes the global model information \( Inf_b^r \), which contains the actual global model \( M^{r-1} \) and metadata.
    \item \textbf{On-Chain Action:} He calls the \texttt{PublishTask} smart contract function. This posts the \textit{hash} \( h_{Inf_b^r} \) and other task details (like round number \( r \), project ID \( id_p \), and deadline \( D_t \)) to the blockchain. This serves as an immutable, public start-signal for the round.
\end{itemize}
\textbf{Time Complexity:} \( O(1) \). The KDF and Hash operations are on fixed-size metadata (only the model \textit{hash} is computed, not on the model itself in this step, as \(Inf_b^r\)'s hash is computed for the \textit{next} step). The blockchain call is a network latency.

\subsection*{2. Server (S) Sends Global Model (Off-Chain)}
After posting the public task, Bob sends the actual model to Alice securely.

\[
msg_b \leftarrow \text{Enc}((Inf_b^r \parallel Tx_p), K_{i,j}) \text{ }
\]
\[
\sigma_b \leftarrow \text{Sig.Sign}(sk_b, msg_b) \text{ }
\]
\begin{itemize}
    \item \textbf{Encryption:} Bob encrypts the full global model info (\( Inf_b^r \), which contains the model \( M^{r-1} \)) using the single-use model key \( K_{i,j} \) he just derived. This provides \textbf{confidentiality}.
    \item \textbf{Authentication:} He signs the \textit{ciphertext} (\( msg_b \)) with his private ECDSA key (\( sk_b \)). This provides \textbf{authenticity} (proves it's from Bob) and \textbf{integrity} (proves the ciphertext wasn't tampered with).
    \item \textbf{Transmission:} He sends \( \langle msg_b, \sigma_b \rangle \) to Alice via the off-chain channel.
\end{itemize}
\textbf{Time Complexity:} \( O(L_M) \). The dominant operation is the symmetric encryption (e.g., AES-GCM) of the model, which is linear in the model size \( L_M \). The signature operation is \( O(1) \).

\subsection*{3. Participant (P) Receives \& Trains}
Alice receives the off-chain message, derives the same key, and trains her model.

\[
(CK_{i,j}, K_{i,j}) \leftarrow KDF_S(CK_{i-1,j})
\]
\[
\text{Sig.Ver}(pk_b, msg_b, \sigma_b) \stackrel{?}{=} 1
\]
\[
(Inf_b^r \parallel Tx_p) \leftarrow \text{Dec}(msg_b, K_{i,j})
\]
\[
m^r \leftarrow \text{Train}(M^{r-1}, \text{Data}_{local})
\]
\begin{itemize}
    \item \textbf{Turn the Ratchet:} Alice sees the \(Tx_p\) event on the blockchain, which signals her to "turn her ratchet." She independently applies the \textit{exact same} \( KDF_S \) function to her copy of \( CK_{i-1,j} \) to derive the \textit{identical} keys \( CK_{i,j} \) and \( K_{i,j} \).
    \item \textbf{Verification:} She verifies Bob's signature \( \sigma_b \) using his public key \( pk_b \).
    \item \textbf{Decryption:} If the signature is valid, she uses the derived model key \( K_{i,j} \) to decrypt \( msg_b \) and recover the global model \( M^{r-1} \). She also checks that the \( Tx_p \) data inside the message matches the data on the blockchain.
    \item \textbf{Training:} She trains the received global model on her private local data to produce her new local model \( m^r \).
\end{itemize}
\textbf{Time Complexity:} \( O(L_M) + O(\text{TrainingComplexity}) \). The crypto cost is \( O(L_M) \) (for decryption) plus a few \( O(1) \) operations. The total cost is dominated by the computationally intensive model training step.

\subsection*{4. Participant (P) Publishes Update (On-Chain)}
After training, Alice commits to her new local model on the blockchain.

\[
h_{Inf_a^r} \leftarrow H(Inf_a^r)
\]
\[
Tx_u \leftarrow \{r, h_{Inf_a^r}, id_p, id_t\} \text{ }
\]
\[
BC \leftarrow \text{UpdateModel}(Tx_u)
\]
\begin{itemize}
    \item \textbf{Commitment:} Alice creates her model information \( Inf_a^r \) (which includes her new local model \( m^r \)) and computes its hash \( h_{Inf_a^r} \).
    \item \textbf{On-Chain Action:} She calls the \texttt{UpdateModel} smart contract function, posting her hash \( h_{Inf_a^r} \) to the blockchain. This does \textit{not} reveal her model; it only serves as an immutable, timestamped \textbf{proof of contribution}.
\end{itemize}
\textbf{Time Complexity:} \( O(L_M) \). The dominant cost is hashing the full local model \( m^r \) (which is part of \( Inf_a^r \)), a linear-time operation.

\subsection*{5. Participant (P) Sends Local Model (Off-Chain)}
With her commitment public, Alice now sends her actual model to Bob.

\[
msg_a \leftarrow \text{Enc}((Inf_a^r \parallel Tx_u), K_{i,j}) \text{ }
\]
\[
\sigma_a \leftarrow \text{Sig.Sign}(sk_a, msg_a) \text{ }
\]
\begin{itemize}
    \item \textbf{Encryption:} Alice encrypts her full local model information \( Inf_a^r \) using the \textbf{same model key \( K_{i,j} \)} for this round.
    \item \textbf{Authentication:} She signs the ciphertext \( msg_a \) with her private ECDSA key \( sk_a \).
    \item \textbf{Transmission:} She sends \( \langle msg_a, \sigma_a \rangle \) to Bob via the off-chain channel.
\end{itemize}
\textbf{Time Complexity:} \( O(L_M) \). Dominated by the linear cost of encrypting the model.

\subsection*{6. Server (S) Receives \& Aggregates}
Bob receives Alice's model, verifies it against the blockchain, and aggregates it.

\[
\text{Sig.Ver}(pk_a, msg_a, \sigma_a) \stackrel{?}{=} 1 \text{ }
\]
\[
(Inf_a^r \parallel Tx_u) \leftarrow \text{Dec}(msg_a, K_{i,j}) \text{ }
\]
\[
H(Inf_a^r) \stackrel{?}{=} h_{Inf_a^r} \text{ (from on-chain } Tx_u)
\]
\[
M^r \leftarrow \text{Aggregate}(m^r, \dots)
\]
\[
Tx_f \leftarrow \{r, sc, h(M^r), id_p, id_t, T\} \text{ }
\]
\[
BC \leftarrow \text{FeedbackModel}(Tx_f)
\]
\begin{itemize}
    \item \textbf{Verification:} Bob verifies Alice's signature \( \sigma_a \).
    \item \textbf{Decryption:} He uses the model key \( K_{i,j} \) (which he derived in Step 1) to decrypt \( msg_a \) and recover her local model \( m^r \).
    \item \textbf{Integrity Check (Crucial):} Bob re-hashes the \( Inf_a^r \) he just decrypted. He then reads the hash \( h_{Inf_a^r} \) that Alice posted to the blockchain in Step 4. If the two hashes match, he has cryptographic proof that the model he received off-chain is the \textit{exact same one} she publicly committed to. This prevents tampering and replay attacks.
    \item \textbf{Aggregation:} If valid, he aggregates her model (e.g., using Federated Averaging \( M^r = \sum_{i=1}^{N} \frac{w_i}{N} m_i^r \)).
    \item \textbf{Feedback (On-Chain):} Bob calls the \texttt{FeedbackModel} smart contract function. This posts a score \( sc \) for Alice (used by the reputation mechanism) and the hash of the \textit{new} global model \( h(M^r) \). This loop then repeats, with \( M^r \) becoming the new \( M^{r-1} \) for the next round.
\end{itemize}
\textbf{Time Complexity:} \( O(L_M) \). This is a sum of linear-time operations: decryption \( O(L_M) \), hashing \( O(L_M) \), and aggregation (a linear pass over model weights) \( O(L_M) \).

\newpage

\section*{Phase 4: Asymmetric Ratchet (Key Update)}

This is a special, server-initiated phase designed to provide \textbf{post-compromise security}. If a participant's key state is compromised, this process establishes a completely new, cryptographically independent root key \( RK_{j+1} \). This renders any previously stolen keys useless for decrypting future messages . The server can trigger this after a certain number of rounds (when the symmetric ratchet threshold \(L_j\) is reached) or if a compromise is suspected .

This phase temporarily modifies the `Send and Receive Model Loop` for a single round to exchange new public keys, as described in Sec 4.2.4 of the paper.



\subsection*{1. Server (S) Initiates Ratchet}
Bob (the server) decides to update the root key. He generates new public keys and piggybacks them onto the normal `PublishTask` message for the current round.

\begin{itemize}
    \item \textbf{Generate New Keys (Local):} Bob creates a fresh set of KEM and ECDH key pairs. These will be the foundation for the new root key.
    \[ (kpk'_b, ksk'_b) \leftarrow \text{KEM.Gen()} \]
    \[ (epk'_b, esk'_b) \leftarrow \text{KE.Gen()} \]
    
    \item \textbf{Publish Task (On-Chain) - Modified Step 1:} Bob creates a modified `PublishTask` transaction. The information payload now includes his new public keys, and the transaction on-chain includes a commitment (hash) to these new keys .
    \[ Inf_b^r \leftarrow \langle r, kpk'_b, epk'_b, M^{r-1}, id_p, id_t, D_t \rangle \]
    \[ Tx_p \leftarrow \{ r, H(Inf_b^r), H(kpk'_b \parallel epk'_b), id_p, id_t, D_t \} \]
    \[ BC \leftarrow \text{PublishTask}(Tx_p) \]
    
    \item \textbf{Send Keys (Off-Chain) - Modified Step 2:} This is the crucial step. Bob encrypts this new information payload using the \textbf{last key from the old ratchet}, \( K_{i,j} \). This ensures that only the intended participant (who can derive \( K_{i,j} \)) can receive the new public keys.
    \[ K_{i,j} \leftarrow KDF_S(CK_{i-1,j}) \]
    \[ msg_b \leftarrow \text{Enc}((Inf_b^r \parallel Tx_p), K_{i,j}) \]
    \[ \sigma_b \leftarrow \text{Sig.Sign}(sk_b, msg_b) \]
\end{itemize}
\textbf{Time Complexity:} \( O(L_M) + O(1) \). The complexity is the standard linear time for encryption (\( O(L_M) \)) plus the constant time for generating the new key pairs (\( O(1) \)).

\subsection*{2. Participant (P) Responds to Ratchet}
Alice sees the modified `PublishTask` transaction, receives the off-chain message, and proceeds to establish the new shared key.

\begin{itemize}
    \item \textbf{Receive \& Decrypt (Local) - Modified Step 3:} Alice turns her old ratchet one last time to derive the final key \( K_{i,j} \). She uses this key to decrypt Bob's message.
    \[ K_{i,j} \leftarrow KDF_S(CK_{i-1,j}) \]
    \[ (Inf_b^r \parallel Tx_p) \leftarrow \text{Dec}(msg_b, K_{i,j}) \]
    She verifies the signature and checks that the hash of the new keys from the message matches the hash in the on-chain transaction. She then extracts Bob's new public keys \( kpk'_b \) and \( epk'_b \).
    
    \item \textbf{Generate New Root Key (Local):} Alice now performs the same steps as in the initial session establishment (Phase 2), but with Bob's new keys.
    \begin{enumerate}
        \item She generates her own new ephemeral ECDH key pair: \( (epk'_a, esk'_a) \leftarrow \text{KE.Gen()} \).
        \item She computes the new hybrid shared secrets:
        \[ SS_e \leftarrow \text{KE.Derive}(esk'_a, epk'_b) \]
        \[ (ct', SS_k) \leftarrow \text{KEM.Encap}(kpk'_b) \]
        \item She derives the new root key, \( RK_{j+1} \), and from it, the first set of keys for the new ratchet chain.
        \[ RK_{j+1} \leftarrow KDF_A(SS_e \parallel SS_k, 0x00) \]
        \[ (CK_{1,j+1}, K_{1,j+1}) \leftarrow KDF_S(RK_{j+1}) \]
    \end{enumerate}
    
    \item \textbf{Publish Update (On-Chain) - Modified Step 4:} After training her local model \( m^r \), her on-chain update is also modified to include commitments to her new key material.
    \[ m^r \leftarrow \text{Train}(M^{r-1}, \text{Data}_{local}) \]
    \[ Inf_a^r \leftarrow \langle r, ct', epk'_a, m^r, id_p, id_t \rangle \]
    \[ Tx_u \leftarrow \{ r, H(Inf_a^r), H(ct' \parallel epk'_a), id_p, id_t \} \]
    \[ BC \leftarrow \text{UpdateModel}(Tx_u) \]
    
    \item \textbf{Send Model (Off-Chain) - Modified Step 5:} This is the second crucial step. Alice encrypts her model and response not with the old key, but with the \textbf{first key from the new ratchet}, \( K_{1,j+1} \).
    \[ msg_a \leftarrow \text{Enc}((Inf_a^r \parallel Tx_u), K_{1,j+1}) \]
    \[ \sigma_a \leftarrow \text{Sig.Sign}(sk_a, msg_a) \]
\end{itemize}
\textbf{Time Complexity:} \( O(L_M) + O(\text{TrainingComplexity}) \). The cryptographic cost is linear in the model size (\(O(L_M)\)), but the overall cost is dominated by the model training.

\subsection*{3. Server (S) Completes Ratchet}
Bob receives Alice's response and now has all the information needed to derive the new root key himself and decrypt her message.

\begin{itemize}
    \item \textbf{Receive \& Generate New Root Key (Local):} Bob sees Alice's modified \( Tx_u \) on-chain and receives her off-chain message \( msg_a \). He verifies her signature and checks the on-chain hash commitment. He then extracts her new public key \( epk'_a \) and the new ciphertext \( ct' \).
    
    He now performs the same derivation as Alice to arrive at the identical new root key and first model key.
    \begin{enumerate}
        \item He computes the new hybrid shared secrets using his new \textit{private} keys:
        \[ SS_e \leftarrow \text{KE.Derive}(esk'_b, epk'_a) \]
        \[ SS_k \leftarrow \text{KEM.Decap}(ksk'_b, ct') \]
        \item He derives the new root key and the first model key:
        \[ RK_{j+1} \leftarrow KDF_A(SS_e \parallel SS_k, 0x00) \]
        \[ (CK_{1,j+1}, K_{1,j+1}) \leftarrow KDF_S(RK_{j+1}) \]
    \end{enumerate}
    
    \item \textbf{Process Model (Local) - Modified Step 6:} Now that he has successfully derived the new model key \( K_{1,j+1} \), he can decrypt Alice's message.
    \[ (Inf_a^r \parallel Tx_u) \leftarrow \text{Dec}(msg_a, K_{1,j+1}) \]
    He verifies the hashes, aggregates her local model \( m^r \), and sends the final feedback transaction to the blockchain.
    \[ M^r \leftarrow \text{Aggregate}(m^r, \dots) \]
    \[ BC \leftarrow \text{FeedbackModel}(Tx_f) \]
\end{itemize}
\textbf{Time Complexity:} \( O(L_M) \). The cost is dominated by the linear-time operations of decryption, hashing, and model aggregation.

The asymmetric ratchet is now complete. The protocol continues from Phase 3, Step 1, using the new chain key \( CK_{1,j+1} \) as the input for the next round's symmetric ratchet.
\newpage

\section*{Phase 5: Threat-Adaptive Ratcheting Mechanism}

The standard PQBFL protocol relies on a fixed symmetric ratcheting threshold, $L_j$, which determines the frequency of computationally expensive post-quantum (Kyber) asymmetric key exchanges. A static $L_j$ forces a rigid trade-off between security and performance constraints. To dynamically address real-time adversarial behavior while minimizing overhead, we introduce the \textbf{Threat-Adaptive Ratcheting Mechanism}.

\subsection*{5.1 Mathematical Model of ThreatMonitor}

The system monitors five distinct categories of security-relevant events continuously. Let $E = \{e_1, e_2, e_3, e_4, e_5\}$ be the set of monitored event types:
\begin{enumerate}
    \item $e_1$: failed signature verification ($w_1 = 1.0$)
    \item $e_2$: hash mismatch in key or model commitments ($w_2 = 0.9$)
    \item $e_3$: client reputation drop on-chain ($w_3 = 0.6$)
    \item $e_4$: anomalous round-trip timing deviation ($w_4 = 0.4$)
    \item $e_5$: stale ratchet (too many rounds since last symmetric reset) ($w_5 = 0.3$)
\end{enumerate}

To ensure older events contribute less to the current security posture, we introduce an exponential decay function based on the time elapsed since the event occurred, $\Delta \tau_i$:
\[
decay(i) = 2^{-\Delta \tau_i / \lambda}
\]
where $\lambda$ is the half-life constant for threat memory. The composite, dynamically normalized threat level at time $t$ is calculated as:
\[
\Theta(t) = \frac{\sum_{i=1}^{|E|} \left( \sum_{k=1}^{N_i} decay(i_k) \cdot w_i \cdot sev(i_k) \right)}{\sum_{i=1}^{|E|} w_i \cdot \max(sev_i)}
\]
where $N_i$ is the number of occurrences of event type $e_i$, and $sev(i_k) \in [0, 1]$ marks the severity. Thus, $\Theta(t) \in [0, 1]$.

\textbf{Example Calculation:}
Assume $\lambda = 10 \text{ sec}$, and within a 5-second interval, a node encounters 3 signature verification failures ($e_1$ with $\Delta \tau_1 = 5, 2, 0$ sec, $sev = 1.0$) and a timing anomaly ($e_4$ at $\Delta \tau_4 = 1$ sec, $sev = 0.5$). 
\[
decay_{e_1} = 2^{-5/10} + 2^{-2/10} + 2^{-0} \approx 0.707 + 0.871 + 1.0 = 2.578
\]
\[
decay_{e_4} = 2^{-1/10} \approx 0.933
\]
The raw threat score is: $(2.578 \times 1.0 \times 1.0) + (0.933 \times 0.4 \times 0.5) = 2.578 + 0.187 = 2.765$. Normalized against a historical maximum threshold, if the max aggregate was 5.0, then $\Theta(t) = \frac{2.765}{5.0} = 0.553$.

\subsection*{5.2 Adaptive Policy Function}

The threat level $\Theta(t)$ directly determines the ratcheting parameter $L_j$ for the current state. The relationship is governed by a power curve to ensure $L_j$ rapidly collapses to maximum security limits under high threat, but remains optimally efficient under normal conditions:
\[
L_j(\Theta) = \lfloor L_{max} - (L_{max} - L_{min}) \times (\Theta(t))^{\alpha} \rfloor
\]
where $L_{max}$ (e.g., 20) is the maximum symmetric steps allowed (peak performance), $L_{min}$ (e.g., 2) is the minimum steps allowed (peak security), and $\alpha$ determines the sensitivity curve.

\textbf{Example Calculation:} 
Using $L_{max} = 20, L_{min} = 2, \alpha = 2.0$, and the previously calculated $\Theta(t) = 0.553$:
\[
L_j(0.553) = \lfloor 20 - (18) \times (0.553)^2 \rfloor = \lfloor 20 - 18 \times 0.306 \rfloor = \lfloor 20 - 5.508 \rfloor = 14
\]
Thus, the system dynamically tightens key rotation frequency from 20 rounds to 14 rounds explicitly due to detected environmental anomalies.

\subsection*{5.3 Security Proofs}

\textbf{Theorem 1: Forward Secrecy Preservation}
\textit{Proof Sketch:} Forward secrecy rests on the security of the Key Derivation Function ($KDF_S$). Let the symmetric ratchet be a one-way function $KDF_S: \{0,1\}^\lambda \to \{0,1\}^{\lambda} \times \{0,1\}^\lambda$, mapping $CK_{i-1} \mapsto (CK_i, K_i)$. Because $KDF_S$ relies on a cryptographic primitive structurally modeled as a PRF (Pseudo-Random Function), reversing the chain to find $CK_{i-1}$ given $CK_i$ requires breaking the underlying PRF assumption. The dynamic variation of $L_j$ strictly acts as a parameter dictating the quantity of iterations $i$, not the mathematical structure of $KDF_S$. Since $\forall i, P(\text{Invert } KDF_S) \le \epsilon_{PRF}$, altering $L_j$ dynamically preserves identical generic forward secrecy limits against retroactive traversal.

\textbf{Theorem 2: Enhanced Post-Compromise Security (PCS)}
\textit{Proof Sketch:} Post-compromise security recovery occurs precisely when an asymmetric ratchet takes place, deriving a new master root key $RK_{j+1}$. The attacker's window of compromise is bounded by the number of symmetric keys derived between asymmetric updates: $K_{compromise} = \{K_{i,j}\}_{i=c}^{L_j}$, where $c$ is the round of compromise intercept. Under a fixed paradigm, the attacker gains $L_j^{fixed} - c$ keys. Under the threat-adaptive protocol, an active attack instantly elevates $\Theta(t) \to 1.0$, rapidly forcing $L_j^{adaptive} \to L_{min}$. As $L_{min} \ll L_j^{fixed}$, the compromise window strictly shrinks: $|K_{compromise}^{adaptive}| < |K_{compromise}^{fixed}|$. The generic adversarial advantage remains computationally bounded by the Kyber IND-CCA assumption:
\[
Adv_A^{PCS} \le \min \{Adv_{\mathcal{A},\text{Kyber}}^{IND-CCA}, Adv_{\mathcal{A},\text{ECDH}}^{PRF-ODH}\}
\]
However, the frequency of PCS recovery effectively truncates the utility lifespan of any compromised symmetric subset.

\subsection*{5.4 Overhead and Time Complexity Analysis}

The adaptive policy introduces constant-time $O(1)$ operations per round: processing the ThreatMonitor parameters and performing deterministic evaluations.
\begin{itemize}
    \item \textbf{Threat Evaluation Complexity:} The state is monitored continuously. $decay$ factors and sum accumulations over predefined subsets resolve in strictly $O(1)$ fixed overhead ($\approx 1\text{-}10 \, \mu s$). 
    \item \textbf{Kyber Asymmetric Avoidance:} Each avoided Kyber Key Encapsulation Mechanism (KEM) instance saves $O(k^2 \cdot n \log n)$ degree-$n$ polynomial multiplications. For Kyber-768 ($k=3$), this yields a massive latency contraction compared to conservative static thresholds.
\end{itemize}

\newpage

\section*{Phase 6: Side-Channel Hardening and Cryptographic Security}

Standard Python implementations of basic cryptographic checks inherently introduce side-channels, mostly timing oracles, which permit adversarial byte-by-byte key extraction. The upgraded framework eliminates these vectors entirely via systemic memory and timing-invariant primitives.

\subsection*{6.1 Constant-Time Cryptographic Primitives}

A critical vulnerability exists when verifying incoming cryptographic commitments mathematically via native equality operators (e.g., \verb|==|). Python's baseline $A == B$ execution breaks at the exact index $i$ where $A[i] \ne B[i]$.

\textbf{Mathematical Time Complexity of Data Vulnerability:}
Let $T_{comp}(A, B)$ be the time to evaluate equality.
\[
T_{vulnerable}(A, B) = \mathcal{O}(i) \text{ where } i = \min \{k \mid A[k] \ne B[k] \}
\]
An attacker iteratively tests byte signatures, measuring execution latency. When response time increases by $\Delta t$ at attempt $X$ at byte $k$, the attacker proves $X=A[k]$, subsequently moving to $k+1$.

\textbf{Constant-Time Mitigation:}
We replace all naive equality paths utilizing \verb|hmac.compare_digest(A, B)|, forcing structural traversal across the entirety of sequence $L$:
\[
T_{hardened}(A, B) = \mathcal{O}(L) 
\]
\textbf{Example:} When comparing a 32-byte hash commitment $h_{epk_a}$, $T_{hardened}$ rigidly computes 32 cycles regardless of whether the deviation exists at byte 0 or byte 31, normalizing power traces and completely denying timing deltas. 

\subsection*{6.2 Authenticated Encryption (AES-GCM) with Random Nonces}

In earlier iterations, AES-CTR was employed without a Message Authentication Code (MAC), rendering the global model weight vectors $M^r$ malleable. By flipping cipher bits, an adversary could systematically alter the decoded weights, committing \textbf{Model Poisoning}. Furthermore, deterministic round-based nonce generators ($nonce = H(\text{"round"} \parallel r)$) risk catastrophic cross-session key-stream overlap.

\textbf{Upgrade Model:}
The protocol was hardened using AES-256-GCM (Galois/Counter Mode). This secures confidentiality while generating a computationally binding integrity MAC tags natively. Nonces are now cryptographically drawn from true randomness ($\text{nonce}_{12} \leftarrow \text{urandom}(12)$).

\textbf{Security and Probability of Collision:}
Let $N$ be the number of total encryptions during the lifetime of a key. Under a 12-byte (96-bit) true random nonce, the geometric probability of any single collision $P_c(N)$ before rotating the symmetric key is quantified by the Birthday Paradox limit:
\[
P_c(N) = 1 - e^{\frac{-N^2}{2 \cdot 2^{96}}} \approx \frac{N^2}{2^{97}}
\]
\textbf{Example Calculation:}
For a federated learning parameter cycle extending up to $N = 2^{30}$ (approx. 1 billion encryptions), the risk evaluates to:
\[
P_c(2^{30}) \approx \frac{2^{60}}{2^{97}} = 2^{-37}
\]
Since the protocol strictly restricts massive $N$ per key via the symmetric ratchet turn ($K_{i,j} \to K_{i+1,j}$ automatically per round), $N$ remains essentially 1. Thus, $P_c(1) \approx 2^{-97}$, rendering a nonce-collision operationally impossible.

\subsection*{6.3 System-Wide Entropy via Random KDF Salts}

The $KDF_A$ (Asymmetric root derivation) heavily depends on the cryptographic entropy of combined $SS_e$ (ECDH secret) and $SS_k$ (Kyber decapsulated secret). Utilizing static salts ($0x00$) theoretically allows heavily resourced attackers to precompute massive rainbow tables indexed against standardized key inputs if round parameters leak.

\textbf{Upgrade Model:}
Salts are independently randomized upon session synthesis ($\text{salt} \leftarrow \text{urandom}(32)$). Let $S_{static}$ denote the search space of the static structure $\approx \mathcal{O}(1)$ pre-compilation factor, forcing reliance strictly on the $SS_k \parallel SS_e$ bits. Randomizing the salt inflates the baseline computation overhead matrix before targeting secrets:
\[
S_{dynamic} = S_{static} \times 2^{256} = \mathcal{O}(2^{256})
\]
This mathematically guarantees computational infeasibility against massive pre-configured rainbow dictionaries, strictly preserving forward secrecy boundaries while inducing no execution latency penalty ($O(1)$ derivation).

\newpage

\section*{Annexure}

\subsection*{P-256 Curve Parameters}
\[
y^2 = x^3 + a x + b \pmod p
\]
where
\begin{align*}
p &= \texttt{0xffffffff00000001000000000000000000000000ffffffffffffffffffffffff} \\
K &= GF(p) \\
a &= K(\texttt{0xffffffff00000001000000000000000000000000fffffffffffffffffffffffc}) \\
b &= K(\texttt{0x5ac635d8aa3a93e7b3ebbd55769886bc651d06b0cc53b0f63bce3c3e27d2604b}) \\
E &= \text{EllipticCurve}(K, (a, b)) \\
G &= E(\texttt{0x6b17d1f2e12c4247f8bce6e563a440f277037d812deb33a0f4a13945d898c296}, \\
     & \texttt{0x4fe342e2fe1a7f9b8ee7eb4a7c0f9e162bce33576b315ececbb6406837bf51f5}) \\
\text{Order} &= \texttt{0xffffffff00000000ffffffffffffffffbce6faada7179e84f3b9cac2fc632551}
\end{align*}

\end{document}
